\chapter{Metodologia e Desenvolvimento}
\label{cap:metodologia}

O desenvolvimento da arquitetura de controle distribuída para o braço robótico \textbf{EB15} baseia-se no particionamento físico e lógico de tarefas mecatrônicas entre unidades computacionais complementares de baixo custo. O objetivo primário desta metodologia é blindar o laço determinístico em tempo real de acionamento mecânico dos motores contra as instabilidades assíncronas causadas por pilhas de conectividade de rede sem fio, telemetria em alta frequência e varredura de sensores angulares absolutos.

Este capítulo detalha a engenharia de hardware e o dimensionamento elétrico do protótipo, abrangendo o microcontrolador mestre ESP32 S3, o co-processador auxiliar Arduino Uno, os circuitos de potência dos drivers TB6600 e CNC Shield v3.1, e a malha de instrumentação baseada nos encoders magnéticos AS5600. Na esfera de software, descreve-se a interconexão UART de baixo payload, a lógica de sincronização instantânea por trigger de hardware dedicado e as diretrizes do firmware sob o sistema operacional FreeRTOS no nó mestre e sob arquitetura bare-metal de baixo nível no nó escravo. Por fim, expõe-se o delineamento dos testes experimentais empregados para a validação dinâmica do sistema.

\section{Arquitetura e Dimensionamento de Hardware}
\label{sec:dimensionamento_hardware}

A topologia do sistema mecatrônico do braço EB15 fundamenta-se no desacoplamento físico das tarefas. Para viabilizar a manipulação com 6 graus de liberdade sem comprometer a estabilidade do movimento, o hardware foi dividido em dois nós de controle complementares, integrados por barramentos de instrumentação e interfaces de potência apropriadas.

\subsection{Nó Mestre: ESP32 S3 e Atuação de Base}
\label{subsec:no_mestre}

O Nó Mestre é centrado no microcontrolador \textbf{ESP32 S3} (Espressif Systems), montado sobre um módulo de desenvolvimento com chip ESP32-S3-WROOM-1. Esta unidade embarca a CPU Xtensa LX7 dual-core de 32 bits, operando a uma frequência de clock nominal de 240 MHz. O processador conta com uma Unidade de Ponto Flutuante (FPU) integrada por hardware, 512 KB de SRAM (\textit{Static Random Access Memory}) interna, 384 KB de ROM (\textit{Read-Only Memory}) e coprocessador de ultrabaixo consumo (ULP). A conectividade de rede é destacada porque o levantamento de Kowsalyadevi e Umamaheswari documenta o uso do ESP32 em aplicações Wi-Fi, justificando o transceptor integrado com suporte a Wi-Fi 802.11 b/g/n (2,4 GHz) e Bluetooth 5 (LE) no nó mestre~\cite{kowsalyadevi2025comprehensive}.

A escolha do ESP32 S3 justifica-se por seu poder computacional expressivo e baixo custo financeiro, fornecendo recursos ideais para assumir o papel de coordenador global do robô manipulador. No Nó Mestre, são centralizados os cálculos matemáticos complexos de cinemática inversa (Equações \ref{eq:wrist_center} e \ref{eq:wrist_orientation}) e de interpolação geométrica em Curva-S de sete segmentos (Fases 1 a 7 descritas no Referencial Teórico).

Adicionalmente, o ESP32 S3 assume o acionamento direto das três juntas inferiores da base do braço robótico, responsáveis por suportar a carga inercial primária e a estrutura física oscilante do robô:
\begin{itemize}
    \item \textbf{Junta 1 (J1 - Rotação da Base):} Movimentada por um motor de passo bipolar NEMA 23 de alto torque (15\,kgf{\cdot}cm), associado a uma caixa de redução planetária impressa em 3D com relação $R_1 = 5:1$.
    \item \textbf{Junta 2 (J2 - Ombro):} Acionada por um motor de passo NEMA 23 de 19\,kgf{\cdot}cm acoplado a uma engrenagem cicloidal estrutural com redução $R_2 = 10:1$, projetada para vencer o momento fletor dinâmico causado pela extensão dos elos estruturais.
    \item \textbf{Junta 3 (J3 - Cotovelo):} Movida por um motor de passo NEMA 17 de alto torque (4{,}8\,kgf{\cdot}cm) associado a uma redução por correias e polias de $R_3 = 4:1$.
\end{itemize}

Devido às altas correntes exigidas pelos motores NEMA 23 e NEMA 17 da base do manipulador, o ESP32 S3 interliga-se a três drivers industriais robustos \textbf{TB6600}. O driver TB6600 utiliza modulação por largura de pulso (PWM --- do inglês, \textit{Pulse Width Modulation}) por corrente constante bipolar para controlar a corrente aplicada aos enrolamentos dos motores. Ele suporta tensões de entrada de até 42\,V e correntes ajustáveis de até 4{,}0\,A, isoladas eletricamente por optoacopladores internos de alta velocidade nas linhas de entrada digital de sinal (Pulso, Direção e Habilitação). A tensão de barramento de potência de alimentação dos drivers TB6600 é mantida estável em 24\,V por meio de uma fonte chaveada de 10\,A. Para atenuar ressonâncias mecânicas locais, as chaves DIP (\textit{Dual In-line Package}) dos drivers TB6600 foram configuradas no modo de micropasso de 1/16 (equivalente a 3200 passos por revolução completa do motor).

\subsection{Nó Escravo: Arduino Uno e Atuação do Punho}
\label{subsec:no_escravo}

O Nó Escravo atua estritamente como um processador executor e aliviador de carga computacional, baseado no microcontrolador monolítico de 8\,bits \textbf{AVR ATmega328P} operando em clock estável de 16\,MHz no módulo \textbf{Arduino Uno}. Esta placa possui 32\,KB de memória flash para armazenamento de código, 2\,KB de SRAM e 1\,KB de EEPROM. A ausência de sistemas operacionais complexos ou pilhas de rede busca reproduzir, em escala de baixo custo, o comportamento temporal previsível exigido por aplicações robóticas em tempo real, como discutido por Økter e Sanfilippo em sistemas hápticos integrados~\cite{okter2025ros}.

O Nó Escravo é acoplado diretamente a uma placa de expansão de barramento comercial \textbf{CNC Shield v3.1} (do inglês, \textit{Computer Numerical Control}), projetada especificamente para controle de eixos múltiplos. Esta placa de expansão simplifica o layout mecatrônico ao encapsular a pinagem de potência e controle lógico em um encapsulamento compacto. O Arduino Uno gerencia de forma direta as três juntas do punho e efetuador do manipulador EB15:
\begin{itemize}
    \item \textbf{Junta 4 (J4 - Rotação do Punho):} Acionada por um motor de passo NEMA 17 de 3{,}2\,kgf{\cdot}cm acoplado a uma redução helicoidal de $R_4 = 3:1$.
    \item \textbf{Junta 5 (J5 - Flexão do Punho):} Movimentada por um motor de passo compacto NEMA 14 de 1{,}5\,kgf{\cdot}cm e redução direta de $R_5 = 2:1$.
    \item \textbf{Junta 6 (J6 - Rotação e Atuação da Pinça):} Controlada por um micro-motor de passo NEMA 11 de 0{,}8\,kgf{\cdot}cm para manipulação e fechamento das garras.
\end{itemize}

O acionamento desses atuadores de menor porte é realizado por três drivers modulares de passo \textbf{A4988} (ou DRV8825), montados diretamente nos soquetes dedicados da CNC Shield v3.1. O A4988 opera com tensões de barramento lógico-mecânico de 12\,V, fornecidas por uma fonte auxiliar chaveada de 5\,A. O modo de micropasso é configurado via pontes de jumpers elétricos por hardware localizados abaixo do driver na placa CNC Shield, ajustados no limite de 1/16 micropasso. A citação a Dong contextualiza a necessidade de reduzir atrasos e irregularidades de transmissão no acionamento do punho, preservando um comportamento suave e atenuado do subsistema mecatrônico~\cite{dong2022comparative}.

\subsection{Sistema de Telemetria e Realimentação Absoluta}
\label{subsec:telemetria_sensores}

Para mitigar a perda de passo mecânico e monitorar desvios posicionais decorrentes da flexibilidade e folgas inerentes à manufatura aditiva 3D do braço EB15, implementou-se uma malha fechada local de realimentação baseada em sensores magnéticos absolutos \textbf{AS5600} (ams OSRAM) instalados individualmente nas 6 juntas articuladas. Essa escolha segue a modelagem de Tang e Yu para controle de motores de passo com encoder magnético~\cite{tang2022research}.

O circuito integrado AS5600 baseia-se em uma matriz interna de sensores de efeito Hall de alta sensibilidade que detecta a orientação espacial do vetor campo magnético gerado por um pequeno ímã de neodímio diametralmente magnetizado, o qual deve ser colado exatamente no centro geométrico do eixo de saída rotativo de cada elo articulado. A caracterização física de Tang e Yu sustenta a descrição da conversão analógico-digital de 12\,bits e da saída angular absoluta em 4096 divisões discretas por volta de 360 graus, correspondente a uma resolução angular real de $0{,}088^{\circ}$ por divisão~\cite{tang2022research}.

A leitura dessa telemetria ocorre via protocolo serial síncrono bi-direcional \textbf{I\textsuperscript{2}C} operando em modo de alta velocidade (\textit{Fast Mode} a 400\,kHz). Uma decisão de projeto relevante no braço EB15 consiste na fixação física dos sensores AS5600 diretamente nos eixos de rotação de saída das juntas estruturais impressas em 3D, após a etapa física de redução mecânica, e não nas partes traseiras dos eixos lineares internos dos motores de passo. Essa disposição aplica ao protótipo a preocupação de Dombre e Khalil com erro geométrico e desempenho de manipuladores, pois permite que a realimentação meça a rotação real do elo e capture folgas mecânicas (\textit{backlash}) acumuladas nas transmissões~\cite{dombre2007robot}.

A distribuição de varredura I\textsuperscript{2}C foi dividida de forma a evitar gargalos de barramento:
\begin{itemize}
    \item \textbf{Canais I\textsuperscript{2}C do Mestre (ESP32 S3):} O ESP32 S3 utiliza seus dois canais físicos independentes de hardware configuráveis por software em seus pinos GPIO para realizar a varredura síncrona dos sensores das juntas inferiores da base (J1, J2 e J3). A leitura é conduzida de forma paralela no Core 1 do ESP32 S3 na taxa fixa de amostragem do laço fechado (200\,Hz).
    \item \textbf{Canal I\textsuperscript{2}C do Escravo (Arduino Uno):} O Arduino Uno adota uma solução de \textit{Software I\textsuperscript{2}C} (\textit{bit-banging}) nos pinos analógicos \texttt{A2} e \texttt{A3}, operando a aproximadamente 200\,kHz. Esta abordagem elimina o conflito de barramento com a biblioteca \texttt{SoftwareSerial}, que reserva os pinos \texttt{A4} (SDA) e \texttt{A5} (SCL) do canal I\textsuperscript{2}C de hardware nativo para varrer e processar localmente a leitura posicional dos encoders magnéticos das juntas superiores do punho (J4, J5 e J6). A amostragem local ocorre a 200\,Hz, sincronizada por interrupções periódicas de temporizadores internos.
\end{itemize}

A separação física dos barramentos I\textsuperscript{2}C impede que ruídos eletromagnéticos propagados ao longo do braço ou colisões de pacotes gerem atrasos nos laços de controle das placas correspondentes.

\section{Interconexão e Protocolo de Sincronização Embarcada}
\label{sec:protocolo_sincronizacao}

O acoplamento elétrico e lógico entre o Nó Mestre e o Nó Escravo representa o ponto crítico para assegurar o paralelismo e a partida simultânea do movimento nas 6 juntas do manipulador robótico.

\subsection{Conectividade e Adequação de Nível Elétrico}
\label{subsec:nivel_eletrico}

A interconexão de controle básico entre o ESP32 S3 e o Arduino Uno é efetuada por meio de uma interface física de comunicação assíncrona cruzada tipo \textbf{UART}. A fiação elétrica interliga o pino de transmissão física TX (GPIO 17) do ESP32 S3 diretamente ao pino de recepção física RX (GPIO 0) do Arduino Uno, e o pino de transmissão TX (GPIO 1) do Arduino Uno ao pino RX (GPIO 18) do ESP32 S3, além de conectar diretamente as linhas de referência de potencial de terra (\textit{Signal GND}) de ambos os circuitos eletrônicos.

Devido à incompatibilidade de tensão lógica entre as duas plataformas mecatrônicas, uma proteção de barramento de alta velocidade torna-se necessária. O ESP32 S3 opera nativamente em lógica CMOS (\textit{Complementary Metal-Oxide-Semiconductor}) de 3{,}3\,V, de forma que seus pinos de entrada e saída digitais toleram tensões máximas equivalentes à sua alimentação lógica de 3{,}3\,V. Em contrapartida, o Arduino Uno (ATmega328P) opera sob lógica TTL (\textit{Transistor-Transistor Logic}) estável de 5{,}0\,V, emitindo sinais de nível lógico alto próximos a 5\,V em seu pino de transmissão TX. A injeção direta de um sinal de 5\,V na linha RX do ESP32 S3 causaria degradação térmica rápida por sobretensão, ativando os diodos de proteção interna contra descarga eletrostática (ESD --- do inglês, \textit{ElectroStatic Discharge}) e levando à queima permanente da CPU mestre.

Para solucionar essa fragilidade física, implementou-se um circuito divisor de tensão resistivo simétrico de alta resposta em frequência na linha de transmissão que parte do TX (5\,V) do Arduino Uno em direção ao RX (3{,}3\,V) do ESP32 S3. O divisor de tensão utiliza um resistor de $1\text{ k}\Omega$ conectado em série na linha de sinal e um resistor de $2\text{ k}\Omega$ conectando a linha do pino RX do ESP32 S3 ao terra digital (GND). A tensão resultante no pino do ESP32 S3 é calculada pela lei de Ohm conforme a Equação~\ref{eq:divisor_tensao}:
\begin{equation}
\label{eq:divisor_tensao}
V_{RX\_ESP32} = V_{TX\_ARDUINO} \cdot \left( \frac{R_2}{R_1 + R_2} \right) = 5{,}0 \cdot \left( \frac{2000}{1000 + 2000} \right) = 3{,}33\text{ V}.
\end{equation}

A Equação~\ref{eq:divisor_tensao} atesta que a tensão é mantida dentro das especificações de segurança lógica do ESP32 S3. Na linha oposta, que transmite o sinal de saída TX (3{,}3\,V) do ESP32 S3 para a entrada RX (5\,V) do Arduino Uno, não há necessidade de adaptação de hardware, uma vez que a tensão limite mínima exigida pelo microcontrolador ATmega328P para registrar com segurança um nível lógico alto digital ($V_{IH}$) é de 3{,}0\,V (equivalente a $0{,}6 \times V_{CC}$), limite atendido pela amplitude nativa do ESP32 S3.

\subsection{Protocolo Serial Compacto Customizado}
\label{subsec:protocolo_customizado}

Para evitar que a latência e o overhead de formatação do canal serial UART degradem o tempo de resposta dinâmico global do sistema, o projeto adota um protocolo customizado estruturado de dados compactados com reduzido payload. Em vez de enviar coordenadas complexas de trajetória em formato textual XML (\textit{eXtensible Markup Language}) ou JSON (\textit{JavaScript Object Notation}), o mestre serializa os dados dinâmicos das juntas superiores (J4, J5 e J6) em frames binários de tamanho fixo. Essa decisão combina as recomendações de Tanvir et al. sobre otimização de payload com a comparação de Dong sobre eficiência de protocolos seriais, justificando o uso de passos, direções pré-calculadas e \textit{checksum} de execução linear rápida no Arduino Uno~\cite{tanvir2024optimizing, dong2022comparative}.

\begin{table}[h!]
\centering
\caption{Estrutura do Frame Binário Customizado para UART (10 Bytes)}
\label{tab:frame_serial}
\begin{tabular}{|c|c|c|c|}
\hline
\textbf{Byte Posição} & \textbf{Campo Técnico} & \textbf{Tipo de Dado} & \textbf{Descrição do Conteúdo} \\ \hline
0 & Preâmbulo Sync & \texttt{uint8\_t} & Padrão fixo de sincronismo elétrico (\texttt{0xAA}) \\ \hline
1--2 & Passos Junta J4 & \texttt{int16\_t} & Número de passos discretos para a Junta J4 (com sinal) \\ \hline
3--4 & Passos Junta J5 & \texttt{int16\_t} & Número de passos discretos para a Junta J5 (com sinal) \\ \hline
5--6 & Passos Junta J6 & \texttt{int16\_t} & Número de passos discretos para a Junta J6 (com sinal) \\ \hline
7 & Parâmetro Velocidade & \texttt{uint8\_t} & Índice da frequência máxima mapeada do Timer \\ \hline
8 & Parâmetro Aceleração & \texttt{uint8\_t} & Rampa de aceleração dinâmica sintonizada \\ \hline
9 & Soma Modular Checksum & \texttt{uint8\_t} & Verificação de paridade (XOR --- do inglês, \textit{eXclusive OR} --- cumulativo de bytes 0 a 8) \\ \hline
\end{tabular}
\end{table}

A transmissão do frame binário compactado leva à seguinte vantagem temporal: à taxa de transmissão física estabelecida de 115200\,bps, cada bit físico transmitido ocupa um intervalo fixo de $8{,}68\ \mu\text{s}$. Considerando que cada byte enviado necessita de 10\,bits de dados de transmissão real (1 start bit, 8 bits de dados e 1 stop bit), a transmissão total do frame físico completo de 10\,bytes é executada em apenas $868\ \mu\text{s}$ ($0{,}868\text{ ms}$). Esta latência sub-milissegundo garante que o canal serial não sature o fluxo de controle em tempo real.

\subsection{Sincronização por Disparo Físico de Hardware}
\label{subsec:trigger_hardware}

Embora a transmissão serial seja otimizada para ocorrer em menos de um milissegundo, a comunicação assíncrona UART apresenta atrasos intrínsecos variáveis decorrentes de latências de processamento interno de software, preenchimento de buffers circulares de recebimento por interrupção e tempo de processamento algorítmico do interpretador da placa receptora. A comparação analítica de Dong é citada para fundamentar essa limitação temporal de protocolos assíncronos, que neste projeto é tratada como uma fonte de \textit{jitter} entre 1 e 3\,milissegundos~\cite{dong2022comparative}.

Se o movimento físico das juntas J1 a J6 dependesse diretamente da recepção assíncrona da mensagem UART pelo Arduino Uno escravo para iniciar o trajeto, haveria um desalinhamento temporal na partida dinâmica dos motores das juntas da base e do punho. A referência a Htun, Hlaing e Hla é usada porque trabalhos de sincronização mestre-escravo em braços robóticos evidenciam que esse tipo de dessincronização compromete a integridade da cinemática cartesiana do efetuador final~\cite{htun2023master}.

Para solucionar esse desvio e obter sincronismo na partida em nível de microssegundo, implementou-se um protocolo de sincronização em duas etapas físicas, integrando o barramento UART a uma linha digital física de disparo por hardware (\textit{hardware trigger}), conforme detalhado a seguir:
\begin{enumerate}
    \item \textbf{Etapa de Configuração e Handshake de Software:}
    Ao planejar uma trajetória no espaço cartesiano, o Nó Mestre ESP32 S3 converte as trajetórias das três juntas superiores do punho em dados de passos, direções e rampas discretas. Em seguida, serializa e transmite o frame binário de 10 bytes para a placa Arduino Uno através do barramento UART. Ao receber o preâmbulo \texttt{0xAA} via recepção serial síncrona por interrupção, o Arduino Uno suspende qualquer varredura secundária, extrai o payload dinâmico, verifica a integridade do \textit{checksum} linear e realiza a pré-configuração física de seus registradores internos de hardware dos temporizadores (\textit{Timer1} e \textit{Timer2}) com os valores recebidos de rampa de aceleração e passos.
    
    Uma vez configurado e pronto para disparar a movimentação em suas três juntas locais, o Arduino Uno transmite um byte simples de prontidão (\texttt{0x06} --- correspondente ao controle ASCII de confirmação \textit{Acknowledge}) via UART de volta para o ESP32 S3. Durante todo esse período de configuração e verificação, o mestre ESP32 S3 permanece no estado lógico de bloqueio, impedido de avançar as juntas J1 a J3 de base.
    
    \item \textbf{Etapa de Disparo Físico de Sincronismo por Hardware:}
    No instante em que o firmware do ESP32 S3 recebe o byte de confirmação \texttt{0x06} enviado pela placa escrava, o Nó Mestre altera o estado lógico de uma linha digital física dedicada GPIO por hardware. Esta linha conecta diretamente um pino de saída digital rápido do ESP32 S3 (GPIO 4) ao pino digital \texttt{9} do microcontrolador ATmega328P no Arduino Uno, configurado como entrada de trigger.
    
    A comutação de estado elétrico (borda de descida do sinal elétrico de 5\,V para 0\,V --- \textit{Falling Edge}) propaga-se ao longo da linha física de cobre em tempo desprezível (frações de nanossegundos). No Arduino Uno, a recepção da borda de descida no pino \texttt{9} é monitorada via \textit{polling} contínuo de sub-microssegundo por software. Essa detecção em baixíssima latência (apenas 180\,ns, equivalentes a 3 ciclos de clock a 16\,MHz) comuta instantaneamente a partida e ativação física dos temporizadores de passo. O uso desse disparo simultâneo aplica os princípios de sincronização determinística discutidos por Nissov, Khedekar e Alexis, alinhando o acionamento do punho ao timer de geração de passos das juntas J1 a J3 no mestre ESP32 S3~\cite{nissov2025simultaneous}.
\end{enumerate}

Essa abordagem de paralelismo determinístico elimina o \textit{jitter} da comunicação UART, garantindo a sincronização espacial síncrona na partida simultânea física de todas as 6 juntas do braço robótico antropomórfico EB15 em patamares sub-milissegundo de precisão mecatrônica.

\section{Arquitetura de Software e Firmware do Sistema Distribuído}
\label{sec:algoritmo_controle}

O projeto lógico do firmware das duas plataformas embarcadas foi estruturado de forma a isolar as atividades assíncronas de rede das rotinas críticas de geração de pulso em tempo real.

\subsection{Firmware Concorrente FreeRTOS no ESP32 S3}
\label{subsec:rtos_mestre}

O firmware do Nó Mestre ESP32 S3 foi desenvolvido em linguagem C/C++ utilizando a estrutura de desenvolvimento oficial do ESP-IDF integrada ao ecossistema Arduino. Para gerenciar a concorrência computacional entre a interface de conectividade e as tarefas críticas de controle de movimento, o sistema adota o sistema operacional de tempo real \textbf{FreeRTOS} para realizar o particionamento de tarefas e fixar sua execução em núcleos de hardware específicos da CPU Xtensa LX7 dual-core.

A distribuição de tarefas do FreeRTOS foi projetada da seguinte forma:
\begin{itemize}
    \item \textbf{Core 1 (Núcleo de Tempo Real Crítico):}
    O núcleo 1 executa as tarefas prioritárias de controle mecânico e realimentação. A tarefa principal \texttt{Task\_BaseControl} opera em alta prioridade (prioridade 10 no escalonador preemptivo) e possui um período de amostragem fixo de 5\,milissegundos ($200\text{ Hz}$). Suas atribuições envolvem a leitura dos encoders absolutos AS5600 das juntas da base (J1, J2 e J3) via barramento I\textsuperscript{2}C local, execução da malha fechada de controle de posição PID discreta, aplicação da modulação não linear por tangente hiperbólica e geração física determinística dos trens de pulsos de passo de alta frequência direcionados aos drivers TB6600.
    
    A tarefa de temporização síncrona é acoplada a um temporizador por hardware interno do ESP32 S3 (\textit{Hardware Timer 0}), garantindo desvios temporais negligenciáveis. Também no Core 1 é executada a tarefa \texttt{Task\_SerialComm} (prioridade 8), responsável por serializar e transmitir os pacotes dinâmicos via UART e aguardar o handshake elétrico do Arduino Uno.
    
    \item \textbf{Core 0 (Núcleo de Serviços Assíncronos e Conectividade):}
    O núcleo 0 executa as rotinas relacionadas às pilhas de protocolo de rede sem fio, que demandam processamento assíncrono e imprevisível. A tarefa \texttt{Task\_WiFi} (prioridade 3) gerencia a estabilidade da interface de rede sem fio nos modos simultâneos \textbf{Station (STA)} e \textbf{Access Point (AP)}. 
    
    A tarefa \texttt{Task\_WebServer} (prioridade 2) hospeda o servidor HTTP (\textit{HyperText Transfer Protocol}) embarcado do robô. Os arquivos da interface gráfica do usuário (HTML, CSS e JavaScript compilados) são armazenados na memória flash do chip por meio do sistema de arquivos LittleFS do ESP32 S3 e são transmitidos síncronamente via requisições web. 
    
    A tarefa \texttt{Task\_RTDE} (prioridade 4) processa de forma contínua conexões por soquete TCP/IP na porta padrão de rede, implementando a lógica clássica do protocolo RTDE para transferência de frames de telemetria em tempo real a 50\,Hz e processamento de alvos cartesianos originados de scripts supervisores.
\end{itemize}

Ao designar as tarefas assíncronas do Wi-Fi e do servidor web exclusivamente para o Core 0, e as tarefas críticas de movimento para o Core 1, elimina-se a interferência do tráfego de dados sobre o laço determinístico de controle mecânico. Essa separação é fundamentada por Leng et al., ao relacionarem desempenho de controle e minimização de \textit{jitter}, e por Bañuelos García et al., ao caracterizarem variações temporais em plataformas ESP32~\cite{leng2014improving, banuelos2026timing}.

\subsection{Interface de Teleoperação e Endpoints RTDE}
\label{subsec:interface_rtde}

A interface gráfica de teleoperação fornece um painel web intuitivo e responsivo compatível com navegadores em computadores e dispositivos móveis. A página web permite:
\begin{itemize}
    \item \textbf{Modo Jog Manual:} Controle individual com passos incrementais programáveis (em graus ou milímetros) para todas as 6 juntas e controle manual de abertura e fechamento do efetuador final (J6).
    \item \textbf{Programação de Pontos:} Gravação interativa de trajetórias baseada na inserção sequencial de posições angulares ou cartesianas em uma tabela dinâmica armazenada em flash, permitindo a reprodução cíclica posterior de trajetórias (\textit{Teach-and-Play}).
    \item \textbf{Configuração de Relações de Redução:} Ajuste dinâmico independente das relações de engrenagens de transmissão mecânica ($R_1$ a $R_6$) para cada uma das juntas. Esta funcionalidade mecatrônica é salva em arquivo JSON local e carregada dinamicamente na memória para corrigir os parâmetros de conversão física de passos por grau no modelo cinemático computado na CPU.
\end{itemize}

O tráfego de telemetria em tempo real entre a interface web e o firmware é unificado por meio do protocolo customizado \textbf{RTDE}, operando sobre canal de soquete TCP dedicado. A opção por pacotes compactados é coerente com a análise de Dong sobre eficiência de transmissão, pois o protocolo serializa o estado angular real das 6 juntas, a corrente estática estimada dos motores, os erros instantâneos de malha fechada e a temperatura interna reportada pelas placas para permitir monitoração em tempo real do manipulador~\cite{dong2022comparative}.

\subsection{Firmware Dedicado Bare-metal no Arduino Uno}
\label{subsec:bare_metal_escravo}

O firmware executado pelo microcontrolador AVR ATmega328P no Nó Escravo foi programado em linguagem C pura sem o uso de qualquer sistema operacional embarcado. A arquitetura de software baseia-se em um laço infinito de baixo nível (\textit{main loop}) focado estritamente no processamento de requisições seriais síncronas de interrupção do buffer UART e leituras analógico-digitais secundárias.

A geração determinística de passos mecânicos com ruído de fase temporal nulo é realizada por meio da programação direta de registradores internos de temporizadores de hardware do ATmega328P:
\begin{itemize}
    \item \textbf{Timer 1 (16 bits):} Configurado em modo CTC (\textit{Clear Timer on Compare Match}) para comandar com precisão de microssegundos as rampas de velocidade e aceleração e a geração periódica de pulsos elétricos de passo para as duas juntas do punho (J4 e J5).
    \item \textbf{Timer 2 (8 bits):} Configurado de forma análoga com prescaler específico de hardware para comutar os pulsos de movimentação e controle estático do efetuador final (J6).
\end{itemize}

No exato instante elétrico em que a placa escrava conclui o handshake UART, ela entra em polling estrito no pino digital \texttt{9}. Assim que o ESP32 S3 comuta o trigger físico de sincronismo por hardware, o Arduino Uno detecta a borda de descida do sinal, comutando instantaneamente a partida e ativação física dos temporizadores de passo. A estratégia combina o princípio de sincronização física simultânea discutido por Nissov, Khedekar e Alexis com a preocupação de Bañuelos García et al. com caracterização temporal de \textit{jitter}, garantindo a partida física dos pulsos em sincronia com o Nó Mestre~\cite{nissov2025simultaneous, banuelos2026timing}.

\subsection{Estratégias de Controle em Malha Fechada e Correção de Posição}
\label{subsec:malha_fechada_metodo}

Para implementar a malha fechada de controle posicional para todas as juntas do braço EB15, o desvio angular instantâneo do elo é calculado periodicamente em cada laço de execução a 200\,Hz. O erro de posicionamento $e(t)$ é definido pela Equação~\ref{eq:erro_instantaneo}:
\begin{equation}
\label{eq:erro_instantaneo}
e(t) = \theta_{des}(t) - \theta_{real}(t).
\end{equation}

Onde $\theta_{des}(t)$ é o ângulo teórico planejado da junta computado no planejamento cartesiano e $\theta_{real}(t)$ é o ângulo físico real medido pelo sensor magnético absoluto AS5600 no elo correspondente. O erro $e(t)$ alimenta o algoritmo discreto do controlador PID linear mapeado na Equação~\ref{eq:pid_model}. As variáveis discretizadas de controle são representadas conforme a Equação~\ref{eq:pid_discreto}:
\begin{equation}
\label{eq:pid_discreto}
u[k] = K_p e[k] + K_i \sum_{j=0}^{k} e[j] \Delta t + K_d \left( \frac{e[k] - e[k-1]}{\Delta t} \right).
\end{equation}

Onde $\Delta t = 5\text{ ms}$ é o tempo discreto de amostragem estabelecido, e $u[k]$ é a variável de atuação eletrotécnica linear correspondente. Para atenuar os efeitos dinâmicos de atrito e inércia do braço impresso em 3D, a ação de velocidade gerada sobre os motores de passo obedece à modulação não linear baseada na função de tangente hiperbólica descrita na Equação~\ref{eq:fm_tangente}:
\begin{equation}
\label{eq:fm_tangente}
f(e[k]) = f_{max} \cdot \tanh(\gamma \cdot e[k]).
\end{equation}

Nesta formulação, o valor linear computado do erro de malha fechada $e[k]$ é modulado em frequência não linear de pulso $f(e[k])$ enviada diretamente ao driver de potência (TB6600 ou CNC Shield). Para desvios posicionais expressivos, a função de tangente hiperbólica satura suavemente no limite estável $f_{max}$, acelerando o elo até a velocidade máxima autorizada sem picos abruptos de inércia. À medida que a junta converge para o ângulo planejado de destino, o método de Zhou et al. justifica o decaimento suave da frequência, que amortece oscilações elásticas do braço impresso em 3D e reduz sobressinais dinâmicos de posicionamento~\cite{zhou2023frequency}.

\section{Procedimento de Validação Experimental}
\label{sec:validacao}

Para ratificar a eficácia da arquitetura distribuída desenvolvida no braço robótico de 6 eixos EB15, planejou-se e implementou-se uma infraestrutura de testes experimentais rigorosos para medir parâmetros físicos e computacionais, que serão apresentados no Capítulo 4:
\begin{enumerate}
    \item \textbf{Mensuração Qualitativa de Jitter de Temporização:}
    Utilização de um osciloscópio digital de armazenamento de canais múltiplos (Tektronix TDS2024C) conectado diretamente aos pinos físicos de pulsos de passos (GPIOs de geração do ESP32 S3 e pinos correspondentes na CNC Shield do Arduino Uno). O teste mede a estabilidade temporal da emissão de frequências sob condições severas de estresse computacional artificial na CPU mestre, por meio da injeção de tráfego de rede concorrente intenso de requisições HTTP e soquetes RTDE simultâneos a 100 Hz.
    
    \item \textbf{Análise Quantitativa de Sincronismo na Partida:}
    Medição sistemática por meio do osciloscópio digital do intervalo temporal entre o envio físico do sinal de trigger digital de sincronismo de hardware (borda de descida) pelo ESP32 S3 e o início real da comutação da primeira onda quadrada de trem de pulso de micropasso emitida pelas portas de potência correspondentes dos drivers da base e do punho.
    
    \item \textbf{Acurácia Posicional Cartesiana em Malha Fechada:}
    Validação da precisão cinemática estática e dinâmica do efetuador final do braço EB15 ao descrever trajetórias lineares e contornos circulares de precisão no espaço físico mecatrônico. Os desvios geométricos reais são monitorados pelos encoders absolutos AS5600 nas juntas e confrontados com os modelos cinemáticos teóricos pré-calculados. O teste visa quantificar a robustez e estabilidade do controle por modulação de tangente hiperbólica no amortecimento estrutural das oscilações residuais mecânicas das peças impressas em 3D.
\end{enumerate}

Desta forma, consolida-se uma metodologia estruturada que une a eletrônica de controle, interconexão lógica e software embarcado robusto para validar o braço manipulador de baixo custo.
